% Source: https://artofproblemsolving.com/community/c3h2780401_2019_mbmt_team_round__montgomery_blair_math_tournament
% Stable identity: aops-topic-2780401-q115
Let $ r $ be the ratio of the leading coefficients of $ P,Q $, and put $ R=P-rQ $, of degree $ d<2019 $. It is nonzero, since proportional $ P,Q $ would give either zero or degree 20190. Write the associated monic polynomial $ H(t)=t^{10}-7t^9+22t^8+\cdots+c_{10}$ as $(t-r)^mK(t)$ with $ K(r)\ne0 $. Then $ F=R^mQ^{10-m}K(\dfrac{P}{Q})$ has degree $ md+(10-m)2019 $, because all factors other than $ R $ retain their leading terms. Equating to 2019 gives $ m(2019-d)=9\cdot2019 $. Thus $ m\ge9 $. The case $ m=10 $ would require $ 10d=2019 $, impossible. Hence $ m=9,d=0 $, and $ H(t)=(t-r)^9(t-s)$. Comparing the first two coefficients gives $ 9r+s=7 $ and $ 36r^2+9rs=22 $. Consequently $ 45r^2-63r+22=0 $, so $(r,s)=(\dfrac{2}{3},1)$ or $(\dfrac{11}{15},\dfrac{2}{5})$. Therefore
\[\boxed{c_{10}=\dfrac{2^9}{3^9}\quad\text{or}\quad\dfrac{2\cdot11^9}{3^9\cdot5^{10}}}.\]
Both occur: take $ Q=x^{2019}$ and $ P=rQ+1 $. Then $ F=(P-rQ)^9(P-sQ)=(r-s)Q+1 $ has degree 2019.
